\documentclass[10pt,a4paper]{amsart}
\usepackage[margin=19mm]{geometry}
\usepackage{amsmath,amssymb,mathtools}
\usepackage[scr=rsfs,cal=euler]{mathalfa}
\usepackage{xfrac}
\usepackage[unicode,hidelinks,bookmarks=false]{hyperref}
\newcommand{\ie}{i.e.\ }
\newcommand{\cf}{cf.\ }
\newcommand{\resp}{respectively\ }
\newcommand{\Subsection}{\subsection}
\providecommand{\nobreakdash}{\nobreak\hbox{-}}
\setlength{\emergencystretch}{2em}
\multlinegap0pt
\usepackage[shortlabels]{enumitem}
\usepackage{amssymb}
\usepackage{mathtools}
\newtheorem{lemma}{Lemma}
\def\ball{{I\kern -.35em B}}
\newcommand{\bx}{\bar x}
\title{Anisotropic Higher-Order Semiregularity of Degenerate Generalized Equations}
\author{Tom\'a\v{s} Roubal}
\date{Source: arXiv:2607.29114v1 (CC BY 4.0)}
\begin{document}
\maketitle

\noindent\emph{Source and licence.} Anisotropic Higher-Order Semiregularity of Degenerate Generalized Equations. Version arXiv:2607.29114v1 (CC BY 4.0), distributed by arXiv under the Creative Commons Attribution 4.0 International licence, http://creativecommons.org/licenses/by/4.0/, as recorded in the arXiv licence metadata for this exact version and shown on the abstract page https://arxiv.org/abs/2607.29114v1. Full licence notice: \texttt{licenses/arxiv-2607.29114-CC-BY-4.0.txt}.\par\smallskip

\noindent\textit{Editorial scope.} Counted unit: the article's lemma lem:p-factor-scaled-error (source0.tex lines 1131--1173). The article's complete original proof of it is delivered with it (source0.tex lines 1175--1293, one \texttt{proof} environment of 119 lines). No mathematics is written, repaired or re-derived: the statement and the proof are the article's own lines, in the article's own notation, and no retained line is substituted or re-flowed.  Retained uncounted as the source-local context the proof needs: lines 1030--1034; lines 1066--1069.  Nothing was omitted from any retained span, so the package carries no omission record.  No retained line contains a citation, so the wrapper carries no bibliography and no bibliography entry.  No retained line contains a figure environment, an image-inclusion command or drawing code, so no image is delivered and the package declares no asset dependency.  Editorial formatting only, in addition: an AMS \texttt{amsart} preamble that restates the article's own declarations and macros used by the retained lines, namely the environments \texttt{lemma} on separate counters, together with the standard packages those lines need.  The retained environments print in this document as Lemma 1 in order of appearance, whereas the article prints the same items with the numbering of its own full document.  The complete native source of the article as distributed in the arXiv e-print for this exact version is bundled unchanged as \texttt{sources/206461/source0.tex}.
	\begin{equation}
		P_i f^{(j)}(\bx)=0
		\quad\text{whenever}\quad 1\le j<i\le p.
		\label{eq:P1}
	\end{equation}

	\begin{equation}
		\Psi_p(h)h=0.
		\label{eq:P4}
	\end{equation}

	\begin{lemma}
		\label{lem:p-factor-scaled-error}
		Assume \eqref{eq:P1}, and let $E\subset X$ be a finite-dimensional
		linear subspace. For every $\varepsilon>0$ there exist $\rho,r>0$ such
		that
		$$
		x_t+\rho t\ball_E\subset U
		\quad\text{for every}\quad t\in(0,r],
		$$
		and the following assertions hold:
		\begin{enumerate}[label={\rm (\roman*)}]
			
			\item For every $t\in(0,r]$ and every $\xi\in\rho t\ball_E$, we have
			\begin{equation}
				\left\|D_t^{-1}
				\bigl(f(x_t+\xi)-f(x_t)-A_t\xi\bigr)
				\right\|_\Sigma
				\le\varepsilon\|\xi\|.
				\label{eq:P7}
			\end{equation}
			
			\item If, in addition, the triangular kernel condition
			\eqref{eq:P4} holds, then, after decreasing $r$ if necessary,
			\begin{equation}
				\left\|D_t^{-1}
				\bigl(f(x_t)-f(\bx)\bigr)
				\right\|_\Sigma
				\le\varepsilon\rho t
				\quad\text{for every}\quad t\in(0,r].
				\label{eq:P10}
			\end{equation}
			
		\end{enumerate}
		Consequently, under \eqref{eq:P4},
		\begin{equation}
			\left\|D_t^{-1}
			\bigl(f(x_t+\xi)-f(\bx)-A_t\xi\bigr)
			\right\|_\Sigma
			\le 2\varepsilon\rho t
			\label{eq:centered-scaled-error}
		\end{equation}
		for every $t\in(0,r]$ and every $\xi\in\rho t\ball_E$.
	\end{lemma}

	\begin{proof}
		Choose $\rho_0,\delta>0$ such that
		$$
		\bx+\delta\ball_X\subset U,
		$$
		and let
		$$
		K_0:=h+\rho_0\ball_E.
		$$
		Since $E$ is finite-dimensional, $K_0$ is compact. After decreasing
		$r_0>0$, we have
		$$
		\bx+tw\in U
		\quad\text{for}\quad
		w\in K_0 \quad\text{and}\quad t\in(0, r_0].
		$$
		
		Taylor's formula, together with \eqref{eq:P1}, gives, uniformly for
		$w\in K_0$,
		$$
		t^{1-i}P_i f'(\bx+tw)|_E
		=
		L_i(w)|_E+o(1)
		\quad\text{as}\quad t\downarrow0,
		$$
		for every $i=1,\ldots,p$. Indeed, for $i\ge2$,
		$$
		\begin{aligned}
			&t^{1-i}P_i f'(\bx+tw)|_E-L_i(w)|_E\\
			&\quad=
			\frac{1}{(i-2)!}
			\int_0^1(1-s)^{i-2}
			P_i\bigl(f^{(i)}(\bx+stw)-f^{(i)}(\bx)\bigr)
			[w]^{i-1}|_E\,ds,
		\end{aligned}
		$$
		while the case $i=1$ follows directly from the continuity of $f'$.
		Consequently,
		$$
		\sup_{w\in K_0}
		\left\|
		D_t^{-1}f'(\bx+tw)|_E-\Psi_p(w)|_E
		\right\|_\Sigma
		\longrightarrow0
		\quad\text{as}\quad  t\downarrow0.
		$$
		
		By continuity of $w\mapsto\Psi_p(w)|_E$, choose
		$\rho\in(0,\rho_0]$ and then $r\in(0,r_0]$ so that
		$$
		\sup_{z\in\rho\ball_E}
		\left\|
		D_t^{-1}f'\bigl(\bx+t(h+z)\bigr)|_E
		-\Psi_p(h)|_E
		\right\|_\Sigma
		\le\varepsilon
		$$
		for every $t\in(0,r]$.
		
		Let $\xi\in\rho t\ball_E$. Since
		$D_t^{-1}A_t=\Psi_p(h)$, the fundamental theorem of calculus gives
		$$
		\begin{aligned}
			D_t^{-1}
			\bigl(f(x_t+\xi)-f(x_t)-A_t\xi\bigr)
			=
			\int_0^1
			\left[
			D_t^{-1}f'(x_t+s\xi)|_E-\Psi_p(h)|_E
			\right]\xi\,ds.
		\end{aligned}
		$$
		Taking norms proves \eqref{eq:P7}.
		
		Suppose now that \eqref{eq:P4} holds. For every
		$i=1,\ldots,p$, Taylor's formula and \eqref{eq:P1} yield
		$$
		P_i\bigl(f(\bx+th)-f(\bx)\bigr)
		=
		\frac{t^i}{(i-1)!}
		\int_0^1(1-s)^{i-1}
		P_i f^{(i)}(\bx+sth)[h]^i\,ds.
		$$
		Condition \eqref{eq:P4} implies
		$$
		P_i f^{(i)}(\bx)[h]^i=0,
		$$
		and hence
		$$
		P_i\bigl(f(\bx+th)-f(\bx)\bigr)=o(t^i).
		$$
		Therefore
		\begin{equation}
			\label{eq:p-factor-center-small-o}
			\left\|D_t^{-1}
			\bigl(f(x_t)-f(\bx)\bigr)
			\right\|_\Sigma
			=
			\sum_{i=1}^p
			t^{1-i}
			\left\|P_i\bigl(f(x_t)-f(\bx)\bigr)\right\|
			=o(t).
		\end{equation}
		Since $\rho>0$ is already fixed, decreasing $r$ proves
		\eqref{eq:P10}.
		
		Finally, \eqref{eq:P7}, \eqref{eq:P10}, and
		$\|\xi\|\le\rho t$ give
		$$
		\begin{aligned}
			\left\|D_t^{-1}
			\bigl(f(x_t+\xi)-f(\bx)-A_t\xi\bigr)
			\right\|_\Sigma\le
			\varepsilon\|\xi\|+\varepsilon\rho t
			\le2\varepsilon\rho t,
		\end{aligned}
		$$
		which proves \eqref{eq:centered-scaled-error}.
	\end{proof}

\end{document}
